\documentclass[a4paper, 11pt, reqno]{amsart}




\usepackage{amsmath, amsthm, amssymb}
\numberwithin{equation}{section}
\usepackage{xspace}
\usepackage{graphicx}
\usepackage{array}
\usepackage{braket}
\usepackage{multicol}
\usepackage{mathtools}
\usepackage{enumerate}
\usepackage{delarray}
\usepackage{mathtools}
\usepackage{faktor} % For quotients
\usepackage{mathrsfs}
\usepackage{tikz-cd}
\usepackage{hyperref}
\usepackage[normalem]{ulem} % For underlining
\usepackage[italicdiff]{physics} % For differentials
\usepackage{bbm} % For indicator
\usepackage{float}
\usepackage{stmaryrd} % for double brackets
\usepackage{aligned-overset}
\usepackage{xcolor}
\usepackage{cite}


\theoremstyle{plain}
\newtheorem{thm}{Theorem}[section]
\newtheorem{lemma}[thm]{Lemma}
\newtheorem{corollary}[thm]{Corollary}
\newtheorem{maintheorem}[thm]{Theorem}
\newtheorem{prop}[thm]{Proposition}
\theoremstyle{remark}
\newtheorem*{claim}{Claim}
\theoremstyle{definition}
\newtheorem{remark}[thm]{Remark}
\newtheorem{exmp}[thm]{Example}
\newtheorem{defn}[thm]{Definition}
\newtheorem*{notation*}{Notation}

\newenvironment{solution}
{\begin{proof}[Solution]}
{\end{proof}}




\DeclarePairedDelimiter{\ceil}{\lceil}{\rceil}


\newcommand{\sub}{\subseteq}
\newcommand{\lcm}{\text{lcm}}
\newcommand{\mc}[1]{\mathcal{#1}}
\newcommand{\mf}[1]{\mathfrak{#1}}
\newcommand{\ms}[1]{\mathscr{#1}}
\newcommand{\sO}{\mathcal{O}}
\newcommand{\cyclic}[1]{\langle#1\rangle}
\newcommand{\units}[1]{#1 ^{\times}}
\newcommand{\la}{\langle}
\newcommand{\ra}{\rangle}
\newcommand{\lr}[1]{\left(#1\right)}
\newcommand{\lrvert}[1]{\left\lvert#1\right\rvert}

\DeclarePairedDelimiterX{\inp}[2]{\langle}{\rangle}{#1, #2}

\newcommand{\usub}{\rightsquigarrow}


\makeatletter
\newcommand*{\bigcdot}{}% Check if undefined
\DeclareRobustCommand*{\bigcdot}{%
  \mathbin{\mathpalette\bigcdot@{}}%
}
\newcommand*{\bigcdot@scalefactor}{.5}
\newcommand*{\bigcdot@widthfactor}{1.15}
\newcommand*{\bigcdot@}[2]{%
  \sbox0{$#1\vcenter{}$}% math axis
  \sbox2{$#1\cdot\m@th$}%
  \hbox to \bigcdot@widthfactor\wd2{%
    \hfil
    \raise\ht0\hbox{%
      \scalebox{\bigcdot@scalefactor}{%
        \lower\ht0\hbox{$#1\bullet\m@th$}%
      }%
    }%
    \hfil
  }%
}
\makeatother

\newcommand{\C}{\mathbb{C}}
\newcommand{\F}{\mathbb{F}}
\newcommand{\E}{\mathbb{E}}
\newcommand{\N}{\mathbb{N}\xspace}
\newcommand{\I}{\mathbb{I}\xspace}
\newcommand{\R}{\mathbb{R}\xspace}
\newcommand{\Z}{\mathbb{Z}\xspace}
\newcommand{\Q}{\mathbb{Q}\xspace}
\newcommand{\G}{\mathbb{G}\xspace}
\newcommand{\B}{\mathbb{B}\xspace}
\newcommand{\M}{\mathbb{M}\xspace}
\newcommand{\T}{\mathbb{T}\xspace}
\renewcommand{\S}{\mathbb{S}\xspace}
\let\P\relax
\DeclareMathOperator*{\P}{\mathbb{P}}

\newcommand{\D}{\mathcal{D}}

\renewcommand{\H}{\mathcal{H}}
\newcommand{\K}{\mathcal{K}}

\newcommand{\rG}{\mathcal{G}}

\DeclareMathOperator{\sign}{sign}
\DeclareMathOperator{\area}{area}
\DeclareMathOperator{\Ind}{Ind}
\DeclareMathOperator{\Rep}{Rep}
\DeclareMathOperator{\VCdim}{VCdim}
\DeclareMathOperator{\GL}{GL}
\DeclareMathOperator*{\argmin}{argmin} % no space, limits underneath in displays
\DeclareMathOperator*{\argmax}{argmax}
\DeclareMathOperator{\Bor}{Bor}
\DeclareMathOperator{\Spec}{Spec}
\DeclareMathOperator{\res}{res}
\DeclareMathOperator{\ord}{ord}
\DeclareMathOperator{\Sym}{Sym}
\DeclareMathOperator{\alb}{alb}
\DeclareMathOperator{\img}{Im}
\DeclareMathOperator{\et}{et}
\DeclareMathOperator{\ck}{coker}
\DeclareMathOperator{\Reg}{Reg}
\DeclareMathOperator{\Cor}{Cor}
\DeclareMathOperator{\Ac}{at}
\DeclareMathOperator{\supp}{supp}
\DeclareMathOperator{\Hom}{Hom}
\DeclareMathOperator{\Pic}{Pic}
\DeclareMathOperator{\Gal}{Gal}
\DeclareMathOperator{\fc}{frac}
\DeclareMathOperator{\Ann}{Ann}
\DeclareMathOperator{\Mod}{Mod}
\DeclareMathOperator{\Cone}{Cone}
\DeclareMathOperator{\FI}{FI}
\DeclareMathOperator{\End}{End}
\DeclareMathOperator{\Alb}{Alb}
\DeclareMathOperator{\Ext}{Ext}
\DeclareMathOperator{\ab}{ab}
\DeclareMathOperator{\Jac}{Jac}
\DeclareMathOperator{\coker}{coker}
\DeclareMathOperator{\fr}{frac}
\DeclareMathOperator{\Int}{Int}
\let\Span\relax
\DeclareMathOperator{\Span}{Span}
\DeclareMathOperator{\Ran}{Ran}
\DeclareMathOperator{\ran}{ran}
\DeclareMathOperator{\ext}{ext}
\DeclareMathOperator{\Prob}{Prob}
\DeclareMathOperator{\graph}{graph}
\DeclareMathOperator{\Aut}{Aut}
\DeclareMathOperator{\Ord}{Ord}
\DeclareMathOperator{\pOrd}{p-Ord}
\DeclareMathOperator{\Vect}{Vect}
\DeclareMathOperator{\Par}{Par}
\DeclareMathOperator{\Proj}{Proj}
\DeclareMathOperator{\cb}{cb}
\DeclareMathOperator{\id}{id}
\DeclareMathOperator{\nor}{nor}
\DeclareMathOperator{\opp}{op}
\DeclareMathOperator{\Ad}{Ad}
\DeclareMathOperator{\Dom}{Dom}
\newcommand{\summ}{\sum\limits}
\newcommand{\thicc}{\bigg}
\newcommand{\eps}{\varepsilon}
\newcommand*\cls[1]{\overline{#1}}
\newcommand{\ind}{\mathbbm{1}}
\DeclareMathOperator{\sgn}{sgn}
\newcommand{\acts}{\curvearrowright}
\let\div\relax
\DeclareMathOperator{\div}{div}
\newcommand{\brac}[1]{\langle#1\rangle}
\newcommand{\pdvN}[3]{\frac{\partial^{#1} #2}{\partial #3^{#1}}}


\newcommand\restr[2]{\ensuremath{\left.#1\right|_{#2}}}





\usepackage{enumitem}
\usepackage{color}
\newlist{steps}{enumerate}{1}
\usepackage[foot]{amsaddr}
\setlist[steps, 1]{label = Step \arabic*:}



\hypersetup{%
  colorlinks=true,%
  linkcolor=blue,%
  citecolor=blue,%
  filecolor=blue,%
  menucolor=blue,%
  urlcolor=blue,%
  pdfnewwindow=true,%
  pdfstartview=FitBH
}   










\title{pp}

\newcommand{\KK}{{\mathcal K}}

\newcommand{\KKK}{{\mathbb K}}
\newcommand{\FF}{{\mathcal F}}
\newcommand{\BB}{{\mathbb{ B}}}

\newcommand{\XX}{{\mathbb{ X}}}
\newcommand{\HH}{{\mathcal H}}
\newcommand{\VV}{{\mathcal V}}

\newcommand{\EE}{{\mathcal E}}

\newcommand{\GG}{{\mathcal G}}

\begin{document}

\title[Properly proximal graph product von Neumann algebras]{Properly proximal graph product von Neumann algebras}





\maketitle

\section{Introduction}

The notion of proper proximality for groups was introduced in \cite{BIP21}, generalizing Ozawa's notion of biexactness from \cite{Oza04}. This notion is a stengthening of non--inneramenability  in the sense of Effros, and has applications to the study of von Neumann algebras. In particular, it is an invariant upto measure equivalence and von Neumann equivalence \cite{IsPeRu19}, and when combined with W$^*$CBAP it yields absence of Cartan subalgbras, following ideas of \cite{OzPo10a}. 

In the paper \cite{DKE24properproximalgroups}, proper proximality was classified for graph products of groups, in the sense of Green \cite{Gr90}. In this short note, this is generalized to the setting of graph products of tracial von Neumann algebras \cite{CaFi17}. Indeed the notion of proper proximality has been extended to von Neumann algebras in \cite{DKEP23}, so this question is natural. First we prove the following result:

 \begin{thm}
     Let $\Gamma$ be a finite graph with $|\VV(\GG)|\geq 3$, and $(M_v,\varphi_v)$ be separable tracial for each $v\in \VV$. If $\Gamma$ is irreducible (in the sense that there does not exists subgraphs $\GG_1,\GG_2$ such that $ \VV(\GG_1)\cup \VV(\GG_2)=\VV(\GG) $ and every $(v,u)\in \EE(\GG) $ for all $v\in\VV(\GG_1),u\in \VV(\GG_2)$), and each $M_v$ has a trace zero unitary, then $M$ is properly proximal.
 \end{thm}  

 A full classification for graph products then follows by combining this statement with the known fact from \cite{DKEP23} that a tensor product of finite von Neumann algebras $N_1\overline{\otimes}N_2$ is properly proximal if and only if both are. Indeed, every graph can be decomposed into combinations of the above and irreducible subgraphs. 

 For the sake of brevity, we will quickly recall and follow the standard notations from \cite{DKEP23}, and assume that the reader is familiar with standard notions and arguments in von Neumann algebra theory.



\section{Proof of main result}
Let $ (\GG,\EE) $ be a finite graph, let $ \{( M_v,\varphi_v ):v\in \VV\} $ be a family of von Neumann algebras with faithful normal states. Consider the graph product von Neumann algebra $(M,\varphi)$.

Let $ \GG_1 $ and $\GG_2$ be two subgraphs, and $M_1$, $M_2$ the corresponding graph product von Neumann algebras, considered as subalgebras of $M$. Let also $N = M_1\cap M_2$ be the graph product of the subgraph $\GG_1\cap \GG_2$. 

Let $e_i$, $i=1,2$ be the projection onto the subspace $L^2(M_i) \subset L^2(M)$. And $e$ be the projection onto $L^2N$.

Consider the the $M$-boundary pieces $ \XX_{M_i} = \overline{  MJMJe_i\BB(L^2M)e_i JMJM}^{\|\cdot \|} $, $  \XX_{N} = \overline{  MJMJe\BB(L^2M)e JMJM}^{\|\cdot \|} $, and their $M$-$M$ $JMJ$-$JMJ$ closure $ \KKK^{\infty,1}_{M_i}(M) $, $\KKK_{N}^{\infty,1}(M)$.

Let $ \iota^\sharp: \BB(L^2M) \to \BB(L^2M)^{\sharp*}_J $ be the canonical embedding.

Let $ q_i = \vee_{u,v\in \mathcal{U}(M)} \iota^{\sharp}(uJvJe_iJv^*Ju^*)$ be the identity of $ (\XX_{M_i})_J^{\sharp*} \subset \BB(L^2M)^{\sharp*}_J$ and $q= \vee_{u,v\in \mathcal{U}(M)} \iota^{\sharp}(uJvJeJv^*Ju^*)$ be the identity of $ (\XX_{N})_J^{\sharp*} $.

We want to show that $\KKK^{1,\infty}_{M_1}(M)\cap \KKK^{1,\infty}_{M_2}(M) = \KKK^{1,\infty}_{N}(M)  $. But notice that $ \KKK^{1,\infty}_{M_i}(M) = (\iota^{\sharp})^{-1}\left( \XX_{M_i}^{\sharp*}  \right) $, it suffices to check that
$$ (\XX_{M_1})_J^{\sharp*}\cap (\XX_{M_2})_J^{\sharp*} = (\XX_{N})_J^{\sharp*}.$$
Finally, since $ (\XX_{M_i})_J^{\sharp*} = q_i \BB(L^2M)^{\sharp *}_Jq_i $, this amounts to show that $ q_1q_2=q_2q_1=q $.

\begin{lemma}
    Let $ \KK $ be a normal representation of $\BB(L^2M)^{\sharp*}_J$ so that $ \BB(L^2M)^{\sharp*}_J\subset \B(\KK) $. Then $\iota^{\sharp}(e_1)q_2\KK =\overline{\iota^{\sharp}( e_1MJMJ e_2)\KK} \subset \overline{\iota^\sharp( MJMJe)\KK} = q\KK$.    
\end{lemma}
\begin{proof}
    Note that \begin{align*}
        e_1MJMJ e_2 =& e_1(M\ominus M_1)J(M\ominus M_1)J e_2+e_1M_2J(M\ominus M_1)J e_2 + e_1MJM_1J e_2\\
        =& e_1(M\ominus M_1)J(M\ominus M_1)J e_2+M_1e_1J(M\ominus M_1)J e_2\\ &+ JM_1Je_1(M\ominus M_1)e_2+ JM_1JM_1e.
    \end{align*}
    So, it suffices to show that the range of those operators belongs to $\iota^\sharp( MJMJe)\KK$.
    
    For this, we first claim that $  (e_1-e)(M\ominus M_1)J(M\ominus M_1)J e_2=0$, or equivalently $e_2(M\ominus M_1)J(M\ominus M_1)J (e_1-e)=0$. Indeed, a vector in $(M\ominus M_1)J(M\ominus M_1)J (e_1-e)L^2M$ is spanned by linear combinations of vectors in $ \HH_{v_1\cdots v_n} $ with $v_1\cdots v_n$ a reduced word in $\VV(\GG)$ such that $ v_j\notin \VV(\GG_2) $ for sum $j$, but this guarantees that the vector is orthogonal to $L^2M_2$.
    
    
    Similarly, we have $ (e_1-e)J(M\ominus M_1)J e_2 =(e_1-e)(M\ominus M_1)e_2 =0$. Therefore, we have $ \iota^{\sharp}(e_1(M\ominus M_1)J(M\ominus M_1)J e_2)\KK \subset \iota^\sharp( e(M\ominus M_1)J(M\ominus M_1)J e_2)\KK \subset \iota^\sharp( MJMJe)\KK  $. Similarly, we have $ M_1e_1J(M\ominus M_1)J e_2, JM_1Je_1(M\ominus M_1)e_2\subset \iota^\sharp( MJMJe)\KK $. Combining those inclusions, we obtain $ \iota^{\sharp}( e_1MJMJ e_2)\KK \subset \iota^\sharp( MJMJe)\KK $.
\end{proof}



\begin{corollary}
    $ [q_1,q_2]=0$ and $q_1q_2=q$.
\end{corollary}
\begin{proof}
    Since $q_2$ commutes with $\iota^{\sharp}(M)$ and $\iota^{\sharp}(JMJ)$, and $q_1 = \vee_{u,v\in \mathcal{U}(M)} \iota^{\sharp}(uJvJe_1Jv^*Ju^*)$, to show $[q_1,q_2] = 0$, it suffices to check that $ [\iota^{\sharp}(e_1),q_2] = 0 $. But by the above lemma, we have $ \iota^{\sharp}(e_1)q_2 = q\iota^{\sharp}(e_1)q_2 = q_2\iota^{\sharp}(e_1)q_2 $. Taking the adjoint, we obtain $ \iota^{\sharp}(e_1)q_2=q_2\iota^{\sharp}(e_1)q_2=  \iota^{\sharp}(e_1)q_2 $.

    Note that as $ [\iota^{\sharp}(e_1),q_2] =0$, we have
    $$ q_1q_2 = q_1 = \vee_{u,v\in \mathcal{U}(M)} \iota^{\sharp}(uJvJe_1Jv^*Ju^*)q_2 = \vee_{u,v\in \mathcal{U}(M)} \iota^{\sharp}(uJvJ) \iota^{\sharp}(e_1)q_2\iota^{\sharp}(Jv^*Ju^*). $$
    Since $ \iota^{\sharp}(e_1)q_2 = q\iota^{\sharp}(e_1)q_2 = q_2\iota^{\sharp}(e_1) = q_2\iota^{\sharp}(e_1)q= q\iota^{\sharp}(e_1)q $, we obtain
    \begin{align*}
        q_1q_2 =& \vee_{u,v\in \mathcal{U}(M)} \iota^{\sharp}(uJvJ) q\iota^{\sharp}(e_1)q\iota^{\sharp}(Jv^*Ju^*)\\ =& q\left( \vee_{u,v\in \mathcal{U}(M)} \iota^{\sharp}(uJvJe_1Jv^*Ju^*) \right)q = qq_1q=q.
    \end{align*} 
\end{proof}


Therefore we immediately obtain the following:

\begin{corollary}\label{corollary for support projections}
    $ \KKK^{1,\infty}_{M_1}(M)\cap \KKK^{1,\infty}_{M_2}(M) = \KKK^{1,\infty}_{N}(M)$. \qed 
\end{corollary}

\begin{lemma}
    Suppose that for is $v\in \VV(\GG)$ $(M_v,\varphi_v)$ is separable and tracial. If $ \{\GG_i\}_{i=1,\cdots,s} $ is a family of subgraphs of $\GG$ and $ M $ is properly proximal relative to $ \Gamma(\GG_i) $ for each $i$, then $ M $ is properly proximal relative to $ \Gamma(\cap_i \GG_i) $. 
\end{lemma}
\begin{proof}
    Suppose that $ M $ is not properly proximal relative to $ \Gamma(\cap_i \GG_i) $, then by \cite[Lemma 8.5]{DKEP23} there is a $ M $-central state $\varphi$ on $ \tilde{\S}_{\XX_{\Gamma(\cap_i \GG_i)}}(M) $ such that $ \varphi|_M $ is normal.

    Let $q_{  \XX_{\Gamma(\cap_i \GG_i)}} $ be the identity of $ (\XX_{\Gamma(\cap_i \GG_i)})_J^{\sharp*} $. If $ \varphi( q_{  \XX_{\Gamma(\cap_i \GG_i)}} ) \neq 0 $, then $ \frac{1}{\varphi( q_{  \XX_{\Gamma(\cap_i \GG_i)}} )} \varphi\circ \text{Ad}(q_{  \XX_{\Gamma(\cap_i \GG_i)}} ) $ is a $ M $-central state $\varphi$ on $ \tilde{\S}_{\XX_{\GG_i}}(M) $ for each $i$ (as $ \text{Ad}(q_{  \XX_{\Gamma(\cap_i \GG_i)}} )\B(L^2M)^{\sharp*}_J \subset (\XX_{\Gamma(\cap_i \GG_i)})_J^{\sharp*} \subset (\XX_{\Gamma( \GG_i)})_J^{\sharp*} \subset \tilde{\S}_{\XX_{\Gamma(\GG_i)}}(M)$), which is a contradiction to the fact that $ M $ is properly proximal relative to $ \Gamma(\GG_i) $.

    If $\varphi( q_{  \XX_{\Gamma(\cap_i \GG_i)}}^{\perp} ) = 1  $, then as $ q_{  \XX_{\Gamma(\cap_i \GG_i)}}^{\perp} = \vee_{i} q_{  \XX_{\Gamma( \GG_i)}}^{\perp} $ (by Corollary \ref{corollary for support projections}), there must be a $i$ such that $ \varphi( q_{  \XX_{\Gamma( \GG_i)}}^{\perp} ) \neq 0 $. Then, since $ \text{Ad}(q_{  \XX_{\Gamma( \GG_i)}}^{\perp})(\tilde{\S}_{  \XX_{\Gamma(\cap_i \GG_i)}}(M)) \subset q_{  \XX_{\Gamma( \GG_i)}}^{\perp}\B(L^2M)^{\sharp*}_Jq_{  \XX_{\Gamma( \GG_i)}}^{\perp} \cap \iota^{\sharp}(JMJ)'\subset  \tilde{\S}_{  \XX_{\Gamma( \GG_i)}}(M)$, we have that $ \frac{1}{\varphi( q_{  \XX_{\Gamma( \GG_i)}}^{\perp} )} \varphi\circ \text{Ad}(q_{  \XX_{\Gamma( \GG_i)}}^{\perp} ) $ is a $ M $-central state $\varphi$ on $ \tilde{\S}_{\XX_{\Gamma(\GG_i)}}(M) $, again contradicting to the fact that $ M $ is properly proximal relative to $ \Gamma(\GG_i) $.
\end{proof}

\begin{lemma}
    Let $ M_1,M_2 $ be two separable tracial von Neumann algebras with a common expected subalgebra $B$. Suppose that there exists unitaries $ u_1,u_2\in \mathcal{U}(M_1) $ and $u_3 \in \mathcal{U}(M_2)$ with $ E_B(u_1)=E_B(u_2)=E_B(u_3)=E_B(u_1^*u_2)=0 $, then the amalgamated free product $ M:= M_1\ast_B M_2 $ is properly proximal relative to $B$.
\end{lemma}
\begin{proof}
 First, we recall the notation $ \lambda_i(e_B) $ is defined as the projection onto the space of vectors in $L^2M $ beginning with $ L^2M_{i+1}^o $ or $L^2B$. Also, we have $ \lambda_i(e_B) \in \S_{B}(M) $ \cite[Prop. 4.16]{toyosawa2025relative}.

    Suppose that $M$ is not properly proximal relative to $B$, then there exists a $M$-central state $ \varphi $ on $\tilde{\S}_B(M)$ that is normal on $M$. We first claim that $\varphi(\iota^\sharp(e_B))=0$. Indeed, suppose $\varphi(\iota^\sharp(e_B))\neq 0$, then $\varphi\circ \text{Ad}(\vee_{u\in \mathcal{U}(M)} \iota^\sharp(ue_Bu^*))$ is a $M$-central positive functional on $\B(L^2M)$ not vanishing on $ Me_BM $ (this is because the projections $\vee_{u\in \mathcal{U}(M)} \iota^\sharp(ue_Bu^*)$ commutes with every unitary in $M$). By \cite[Lemma 3.3]{OP10b} (or \cite[Theorem 2]{BH18amenableabsorption}), this implies that $ M \preceq_M B $ which is impossible (as $\|E_B(x(u_1u_2)^n y)\|_2\xrightarrow{n\to \infty} 0 $ for all $ x,y \in M $). Now, notice that $ \lambda_1(e_B)+\lambda_2(e_B) = 1+e_B$, so $ \varphi(\iota^\sharp(\lambda_1(e_B)+\lambda_2(e_B))) = 1 $. But since
    $$ u_1^*\lambda_1(e_B)u_1 + u_2^*\lambda_1(e_B)u_2\leq \lambda_2(e_B),\quad u_3^*\lambda_2(e_B)u_3\leq \lambda_1(e_B),$$
    we obtain $ 2\varphi( \iota^\sharp(\lambda_1(e_B))) \leq \varphi(\iota^\sharp(\lambda_2(e_B))) \leq \varphi(\iota^\sharp(\lambda_1(e_B)))$ as $\varphi$ is $M$-central. This implies that $ \varphi(\iota^\sharp(\lambda_1(e_B)))=\varphi(\iota^\sharp(\lambda_2(e_B)))=0 $, contradicting to the fact that $ \varphi(\iota^\sharp(\lambda_1(e_B)+\lambda_2(e_B))) = 1$.
\end{proof}

\begin{thm}
    Let $\Gamma$ be a finite graph with $|\VV(\GG)|\geq 3$, and $(M_v,\varphi_v)$ be separable tracial for each $v\in \VV$. If $\Gamma$ is irreducible (in the sense that there does not exists subgraphs $\GG_1,\GG_2$ such that $ \VV(\GG_1)\cup \VV(\GG_2)=\VV(\GG) $ and every $(v,u)\in \EE(\GG) $ for all $v\in\VV(\GG_1),u\in \VV(\GG_2)$), and each $M_v$ has a trace zero unitary, then $M$ is properly proximal.
\end{thm}
\begin{proof}
    For each $v\in \VV(\GG)$, we have
    $$M = \Gamma(\GG/\{v\} )\ast_{\Gamma(\text{link}(v))} \Gamma(\text{star}(v)).$$
    Let $ M_1 =\Gamma(\GG/\{v\} )$, $B = \Gamma(\text{link}(v))$, and $M_2 =\Gamma(\text{star}(v)) $. Since $\Gamma$ is irreducible, either $ M_1 $ must contains two unitaries $ u_1,u_2$ such that $ E_B(u_1)=E_B(u_2)=E_B(u_1^*u_2) =0 $ (otherwise we must have that there is only one vertex $u$ in $\VV(\GG)/\VV(\text{star}(v))$. Now, we must have $(u,u') \in \VV(\GG)$, as otherwise free independency of $M_u $ and $M_{u'}$ would provides the desired unitaries. But this implies that $\GG$ can be decomposed into $ \{u,v\} $ and $ \VV(\GG)/\{u,v\} $, contradicting to that $\GG$ is irreducible). Therefore, by the above Lemma, $ M$ is properly proximal relative to $ \Gamma( \text{link}(v) )$ for each $v$, and thus properly proximal relative to $ \Gamma( \cap_{v\in \VV(\GG)}\text{link}(v) ) =\mathbb{C}$.
\end{proof}

\begin{thebibliography}{DKE24}

\bibitem[BH18]{BH18amenableabsorption}
R\'emi Boutonnet and Cyril Houdayer.
\newblock Amenable absorption in amalgamated free product von {N}eumann algebras.
\newblock {\em Kyoto J. Math.}, 58(3):583--593, 2018.

\bibitem[BIP21]{BIP21}
R\'emi Boutonnet, Adrian Ioana, and Jesse Peterson.
\newblock Properly proximal groups and their von {N}eumann algebras.
\newblock {\em Ann. Sci. \'Ec. Norm. Sup\'er. (4)}, 54(2):445--482, 2021.

\bibitem[CF17]{CaFi17}
Martijn Caspers and Pierre Fima.
\newblock Graph products of operator algebras.
\newblock {\em J. Noncommut. Geom}, 11(1):367--411, 2017.

\bibitem[DEP23]{DKEP23}
Changying Ding, Srivatsav~Kunnawalkam Elayavalli, and Jesse Peterson.
\newblock Properly proximal von {N}eumann algebras.
\newblock {\em Duke Math. J.}, 172(15):2821--2894, 2023.

\bibitem[DKE24]{DKE24properproximalgroups}
Changying Ding and Srivatsav Kunnawalkam~Elayavalli.
\newblock Proper proximality among various families of groups.
\newblock {\em Groups Geom. Dyn.}, 18(3):921--938, 2024.

\bibitem[Gre90]{Gr90}
Elisabeth~Ruth Green.
\newblock {\em Graph products of groups}.
\newblock PhD thesis, University of Leeds, 1990.

\bibitem[IPR19]{IsPeRu19}
Ishan Ishan, Jesse Peterson, and Lauren Ruth.
\newblock Von {Neumann} equivalence and properly proximal groups.
\newblock arXiv:1910.08682, 2019.

\bibitem[OP10a]{OzPo10a}
Narutaka Ozawa and Sorin Popa.
\newblock On a class of {${\rm II}_1$} factors with at most one {C}artan subalgebra.
\newblock {\em Ann. of Math. (2)}, 172(1):713--749, 2010.

\bibitem[OP10b]{OP10b}
Narutaka Ozawa and Sorin Popa.
\newblock On a class of {${\rm II}_1$} factors with at most one {C}artan subalgebra, {II}.
\newblock {\em Amer. J. Math.}, 132(3):841--866, 2010.

\bibitem[Oza04]{Oza04}
Narutaka Ozawa.
\newblock Solid von {N}eumann algebras.
\newblock {\em Acta Math.}, 192(1):111--117, 2004.

\bibitem[TY25]{toyosawa2025relative}
Kai Toyosawa and Zhiyuan Yang.
\newblock On relative biexactness of amalgamated free product von neumann algebras.
\newblock {\em arXiv preprint arXiv:2505.19508}, 2025.

\end{thebibliography}
\end{document}